\documentclass[a4paper,12pt]{ctexart}
\usepackage{xcolor}
\usepackage{graphicx}
\usepackage[top=1.8cm, bottom=1.8cm, left=1.8cm, right=1.8cm]{geometry}
\usepackage{array}
\usepackage{multirow}
\usepackage{hyperref}
\hypersetup{colorlinks=true,linkcolor=blue!70!black}
\linespread{1.35}

\newcommand{\ruby}[2]{#1\textsuperscript{#2}}
\newcommand{\shihai}[2]{%
  \vspace{0.35cm}
  \noindent
  \fcolorbox{orange!60!black}{orange!8}{%
    \begin{minipage}{0.915\textwidth}
    \textbf{\textcolor{orange!70!black}{史海拾遗：}#1} \\[0.1cm]
    {\small #2}
    \end{minipage}%
  }
  \vspace{0.35cm}
}

\title{\textcolor{blue!70!black}{\Huge\textbf{世界古近代史}}}
\author{}
\date{}

\begin{document}
\maketitle
\thispagestyle{empty}

\tableofcontents
\newpage

\section{\textcolor{orange!80!black}{第一课\quad 大河之畔：四大文明古国的诞生}}

在我们的中国古代史课程里，我们看到黄河流域孕育了夏商周三代文明。但中国不是孤独的——差不多同一时期，在北非的尼罗河畔、西亚的两河流域、南亚的印度河流域，另外三个伟大的文明也悄然兴起。

这四个文明有一个惊人的共同点——\textbf{它们都诞生在大河的怀抱里}。这不是巧合。

\subsection*{一、古埃及：尼罗河的赠礼}

希腊历史学家希罗多德说过一句名言："埃及是尼罗河的赠礼。"这话一点不夸张——\textbf{如果没有尼罗河，埃及就是一片沙漠}。

尼罗河每年6月到10月定期泛滥，洪水退去后留下一层肥沃的黑色淤泥。古埃及人不需要施肥，只需要在淤泥上撒种子就能获得丰收。这种几乎"自动"的农业条件让埃及人很早就养活了大批人口，从而有能力建造那些令人瞠目结舌的金字塔。

埃及大约在公元前3100年实现统一。国王被称为\textbf{法老}——他不仅是政治领袖，更是活着的神。埃及人相信法老是太阳神之子，死后灵魂会回到神的世界。这个信念直接催生了埃及最著名的建筑物——\textbf{金字塔}。

金字塔是什么？本质上就是法老的"超级坟墓"。古埃及人相信灵魂不灭，只要尸体保存完好，法老就能在来世继续为王。所以他们发明了\textbf{木乃伊制作技术}——掏空内脏、用香料和盐脱水、用亚麻布层层包裹。现存最大的胡夫金字塔原高146米（相当于一栋40多层的现代大楼），用了约230万块石头，每块平均重2.5吨。在没有起重机的四千五百年前，这是如何做到的？至今仍是一个令人着迷的工程学谜题。

古埃及人还发明了\textbf{象形文字}——一种用图画来表示意思的文字系统。和中国的甲骨文一样，埃及象形文字也是人类最早的文字之一。1799年发现的罗塞塔石碑用三种文字（埃及象形文、埃及草书、古希腊文）刻写了同一段内容，法国学者商\rube{博}{bó}良借此破解了失传千年的埃及文字。

\shihai{为什么四大文明古国都在大河边上？}{
  打开世界地图，你会发现四大文明古国的位置出奇地一致：
  \begin{itemize}
    \item 古埃及 —— 尼罗河
    \item 古巴比伦（两河流域）—— 底格里斯河 + 幼发拉底河
    \item 古印度 —— 印度河 + 恒河
    \item 古中国 —— 黄河 + 长江
  \end{itemize}

  为什么？因为\textbf{农业文明的基本生存逻辑}决定了地理位置的选择。大河定期泛滥带来的泥沙是最天然的肥料；河水提供了稳定的灌溉水源；河流本身是最便捷的运输通道。只有在大河流域，人类才能在技术水平还很低的条件下生产出足够的粮食，养活大量人口，进而发展出复杂的社会组织。

  反过来说——为什么非洲撒哈拉以南、美洲亚马孙雨林在同时期没有出现同等级别的文明？不是因为那里的人"不聪明"，而是因为那里的自然环境（热带雨林土壤贫瘠、缺乏大型可驯化动物等）\textbf{让农业文明的启动门槛太高了}。这不是人的差异，是地理馈赠的差异。理解了这一点，你就不会落入"某些民族天生先进"的偏见——文明的早熟，很大程度上是被地理环境"逼"出来或"养"出来的。
}

\subsection*{二、两河流域：苏美尔、巴比伦与最早的法典}

在今天的伊拉克一带，底格里斯河和幼发拉底河之间有一片形似新月的肥沃土地——\textbf{美索不达米亚}（希腊语"两河之间的土地"）。这里诞生了人类历史上一个极其重要的文明：\textbf{苏美尔}。

苏美尔人在大约公元前3500年就建立了世界上最早的城市。他们发明了\textbf{楔形文字}——不是画图，而是用芦苇杆在软泥板上压出楔形的笔划。因为书写材料是泥板，干燥后极其耐久，大量泥板文书一直保存到今天——我们现在还能读到四千多年前的苏美尔人写的商业合同、学校作业、甚至情诗。

苏美尔之后，两河流域先后兴起了多个帝国。其中最著名的，是\textbf{古巴比伦王国}（约公元前1894—前1595年），以及它的第六代国王\textbf{汉谟拉比}。

汉谟拉比做了一件前无古人的事：他把全国的法律刻在了一根高2.25米的黑色玄武岩石柱上，公之于众——这就是著名的《\textbf{汉谟拉比法典}》。法典有282条，涵盖了婚姻、财产、商业、刑法等各个方面。它最著名的原则是"\textbf{以眼还眼，以牙还牙}"——这是一种"对等报复"原则，在古代其实是一种进步（因为它限制了无限制的复仇：你打掉我一颗牙，我只能打掉你一颗牙，不能要你的命）。

\shihai{汉谟拉比法典 vs 同时期的中国——两种治理思路}{
  汉谟拉比法典大约刻于公元前1754年。同时期的中国是什么状态？大致相当于传说中的夏朝晚期到商朝早期——还没有发现成文法典。

  比较两种治理思路很有意思：
  \begin{itemize}
    \item 汉谟拉比把法律\textbf{刻在石柱上、立在广场中}——理论上每个识字的人都能看到。这是\textbf{法律的公开性}——法律不是君主的秘密武器，而是社会成员都知道的规则。
    \item 中国早期更重视\textbf{礼}——不是用成文的法律条文来约束行为，而是用礼仪、习俗和等级秩序来维护社会。孔子后来总结的儒家思想，本质上延续了这种"靠自觉和教养而非靠惩罚"的治理传统。
  \end{itemize}

  哪种更好？没法简单回答。成文法更透明但更僵硬（"以眼还眼"没有回旋余地）；礼治更灵活但更依赖统治者的品德（君主一坏，全盘皆崩）。这个问题——\textbf{法治还是德治？}——政治学家们争论了几千年，至今没有标准答案。
}

\subsection*{三、古印度：被遗忘的哈拉帕文明}

印度河流域的古代文明，在很长一段时间里是不为人知的。直到1920年代，考古学家在今天的巴基斯坦境内发现了\textbf{哈拉帕}和\textbf{摩亨佐-达罗}两座古城遗址，才揭开了这个文明的面纱。

哈拉帕文明大约存在于公元前2600—前1900年，与埃及古王国和两河流域苏美尔文明差不多同时。它的城市规划和卫生设施令人震惊——摩亨佐-达罗的房屋用标准化的烧砖建造，街道呈网格状布局，而且每家每户都有排水系统和室内厕所，废水通过地下管道排走。这种城市规划水平，当时的世界上只有哈拉帕做得到。

然而，大约在公元前1900年左右，哈拉帕文明突然衰亡了。原因至今争论不休——可能是河流改道、可能是气候变化、可能是一次大规模的地震。哈拉帕的文字尚未被完全破解，所以这个文明的大部分故事仍然是沉默的。

大约在公元前1500年左右，一支自称"雅利安人"（意为"高贵者"）的游牧民族从中亚进入南亚次大陆。他们带来了梵语（印度古典语言）和一套影响深远的社会制度——\textbf{种姓制度}。

种姓制度把社会成员分为四个等级：\textbf{婆罗门}（祭司，最高）、\textbf{刹帝利}（武士和统治者）、\textbf{吠舍}（农民和商人）、\textbf{首陀罗}（仆人，最低）。还有更悲惨的"不可接触者"（贱民），他们的影子被认为是不洁的，连碰都不能碰。一个人的种姓由出生决定，终身不能改变——你出生在哪个种姓，就一辈子在哪层。这套制度在印度延续了三千多年，直到今天在印度农村地区仍有很深的影响。

\subsection*{四、技术与文化}

\textbf{古埃及：}
\begin{itemize}
  \item \textbf{金字塔建造技术}：展示了古埃及人在几何学、工程组织和石材加工方面的惊人能力。
  \item \textbf{太阳历}：埃及人把一年分为365天，分12个月，每月30天，年末加5天节日。这套历法后来被罗马人修订为儒略历，再演变成今天的公历。
  \item \textbf{纸莎草纸}：用尼罗河畔的纸莎草茎制成的书写材料——英语paper（纸）这个单词就来源于papyrus（纸莎草）。
\end{itemize}

\textbf{两河流域：}
\begin{itemize}
  \item \textbf{六十进制}：苏美尔人发明了六十进制，我们今天仍然在用——一小时60分钟、一分钟60秒、一个圆360度，都是苏美尔数学的遗产。
  \item \textbf{车轮和战车}：轮子最早就是在两河流域发明的。看似简单的东西——但在此之前，搬运重物只能靠拖和扛。
  \item \textbf{空中花园}：新巴比伦王国国王尼布甲尼撒二世为思念家乡山水的王妃建造的梯田式花园，被古希腊人列为"世界七大奇迹"之一（是否真实存在尚有争议）。
\end{itemize}

\textbf{古印度：}
\begin{itemize}
  \item \textbf{阿拉伯数字}：其实是印度人发明的（1-9和0），后来由阿拉伯商人传播到欧洲，所以被欧洲人称为"阿拉伯数字"。没有"0"这个概念，数学体系很难发展到今天这种程度。
  \item \textbf{佛教}：公元前6世纪由乔达摩·悉达多（即释迦牟尼——"\ruby{释}{shì}\ruby{迦}{jiā}\ruby{牟}{móu}\ruby{尼}{ní}"意为"\ruby{释}{shì}迦族的圣人"）在印度创立。佛教的核心理念是"众生平等"和"因果轮回"——前者直接挑战了种姓制度。后来佛教在印度衰落了，却在中国、东南亚、日本等地生根发芽。
\end{itemize}

\subsection*{五、本课要点回顾}

\begin{center}
\begin{tabular}{|c|c|p{7cm}|}
\hline
\textbf{文明} & \textbf{兴盛年代} & \textbf{核心标识} \\
\hline
古埃及 & 约前3100—前332年 & 尼罗河灌溉农业；金字塔+木乃伊；象形文字；太阳历 \\
\hline
两河流域 & 约前3500—前539年 & 楔形文字+泥板；汉谟拉比法典；六十进制；车轮 \\
\hline
古印度 & 约前2600—前1900年（哈拉帕）；前1500年后（雅利安） & 城市规划；种姓制度；梵语；印度数字；佛教 \\
\hline
\end{tabular}
\end{center}

\vspace{0.3cm}
{\large\textcolor{orange!80!black}{\textbf{思考与讨论}}}
\begin{enumerate}
  \item 四大文明古国都诞生在大河流域——如果环境太"容易"（食物随处可得）或太"艰苦"（无法种地），文明还会诞生吗？这种"地理决定论"有什么局限性？
  \item 汉谟拉比法典主张"以眼还眼"，你觉得这种"对等复仇"的原则公正吗？和今天的法律精神有什么不同？
  \item 印度的种姓制度延续了三千多年——是什么让一种社会分层制度能够如此长久地存在？你能联想到现代社会中有类似的现象吗（不一定叫"种姓"，但效果类似）？
  \item 阿拉伯数字其实是印度人发明的——历史上很多"命名"其实是"误导"的。你还能想到其他"命名不反映真实起源"的例子吗？
\end{enumerate}

\subsection*{要点总结}
\begin{itemize}
  \item 四大文明古国均诞生于大河流域——大河提供灌溉、肥沃淤泥和交通，是农业文明启动的必要条件
  \item 古埃及以金字塔、木乃伊和象形文字著称，法老兼具政治领袖和神的双重身份
  \item 两河流域的苏美尔人发明了楔形文字和六十进制；汉谟拉比法典是最早的成文法典之一
  \item 古印度的哈拉帕文明有惊人的城市规划，但突然衰亡；随后雅利安人带入的种姓制度影响至今
  \item 文明的早熟在很大程度上受地理环境影响——这不是"民族优劣"的问题，而是地理馈赠的差异
\end{itemize}

\subsection*{讨论题}
\begin{enumerate}
  \item 大河文明的理论说文明诞生需要"适中的环境挑战"——不太容易也不太艰难。你觉得这个理论能解释今天世界的发展差异吗？发达国家大多在温带而非热带——这是地理决定论还是其他原因？
  \item 哈拉帕文明的城市规划（下水道、标准化烧砖）令人震惊——但它在几百年后就衰亡了。一个技术上先进的文明为什么可能比"落后"的文明更脆弱？
\end{enumerate}

\vspace{0.3cm}
\textcolor{blue!70!black}{\textbf{延伸阅读}}
\begin{itemize}
  \item 纪录片《BBC：古代世界》——系统讲述四大古文明的兴衰
  \item 《DK儿童历史百科全书》——有大量精美的古文明复原图和文物照片
\end{itemize}

\newpage

\section{\textcolor{orange!80!black}{第二课\quad 古希腊：民主、哲学与城邦的黄金时代}}

如果说四大古文明代表的是"大河农业帝国"模式，那么古希腊走了一条完全不同的路。

希腊半岛多山、靠海、土地贫瘠，没有尼罗河那样慷慨的大河，也发展不出埃及那样的大一统农业帝国。但恰恰因为\textbf{土地碎片化}，这里形成了一个个独立的\textbf{城邦}——小小的城市国家，每个都有自己的政府、法律和军队。城邦之间互相竞争、互相打仗，但也互相学习。

古希腊最著名的两个城邦，代表了两种完全不同的生活方式：\textbf{斯巴达}和\textbf{雅典}。

\subsection*{一、斯巴达：一个"军营国家"}

斯巴达人是公元前12世纪入侵希腊南部的多利安人后代。他们把当地原住民变成了奴隶（称"希洛人"），但斯巴达的公民人数远远少于希洛人——为了维持统治，斯巴达人过上了全员皆兵的生活。

斯巴达男婴一出生就要经过长老检查：身体不够强壮的直接扔掉。7岁离开家庭进入军营集体生活，12岁开始接受严酷的体能和战斗训练，20岁正式入伍，30岁才能结婚但仍要住在军营，60岁退役。斯巴达人的字典里没有"个人"——\textbf{你的全部存在意义就是为城邦战斗}。斯巴达母亲送儿子上战场时说的不是"注意安全"，而是"带着盾牌回来，或者躺在盾牌上被抬回来"——逃跑的士兵会扔掉沉重的盾牌，所以"带着盾牌回来"的意思是"你没有逃跑"。

斯巴达不是没有文化成就——但和它极端军事化的社会结构相比，文化几乎不值一提。

\subsection*{二、雅典：民主的实验室}

雅典走了一条和斯巴达截然相反的路。

公元前6世纪末，雅典人做了一件在世界政治史上具有划时代意义的事情：他们赶走了\rube{僭}{jiàn}主（通过非法手段夺取政权的独裁者），建立了一套让\textbf{公民直接参与决策}的制度——\textbf{民主}。

雅典的民主是怎么运作的？全体成年男性公民组成\textbf{公民大会}——这是最高权力机构，所有重大决策（宣战、缔约、立法、财政预算）都要经过公民大会投票表决。雅典还有\textbf{五百人议事会}（负责日常政务）和\textbf{陪审法庭}（法官从公民中抽签选出，最多一次抽签选出6000名陪审员）。

抽签是一个很有意思的设计——用随机抽签代替选举。为什么？因为选举往往会选出有钱的、能说会道的人；而抽签意味着\textbf{任何公民都有同等机会参与治理}。这个逻辑虽然朴素，但它触及了民主的一个核心问题：\textbf{普通人有没有能力治理国家？}

\shihai{雅典民主是"真民主"吗？}{
  雅典的民主制度被后世赞美了两千多年，但如果我们用今天的标准去看，它的缺陷是触目惊心的：

  \begin{itemize}
    \item \textbf{只有成年男性公民有政治权利}。雅典总人口约30万，其中公民（含家属）约15万，有权参加公民大会的成年男性公民仅约3万—4万人。妇女、占总人口约三分之一的奴隶、以及住在雅典的外邦人——全部被排除在"民主"之外。
    \item \textbf{直接民主容易变成暴民政治}。最有名的例子是苏格拉底之死——这位古希腊最伟大的哲学家，在公元前399年被雅典公民大会以"不敬神和腐蚀青年"的罪名投票判处死刑。投票结果是280票对220票——一个60票的差距，就处死了一位思想家。
  \end{itemize}

  这引出了一个至今仍然尖锐的问题：\textbf{多数人的决定就一定是对的吗？如果多数人投票决定剥夺少数人的权利，这个决定还算"民主"吗？}现代民主制度给出的答案是：民主不等于"多数人想怎样就怎样"——必须有\textbf{宪法}和\textbf{基本人权}作为所有人（包括少数派）的保护底线。换句话说，\textbf{宪政是对民主的必要约束}。雅典人没有这个意识，苏格拉底就这样死在了民主的投票箱前。
}

\subsection*{三、希波战争：弱小战胜强大的经典案例}

公元前5世纪初，横跨亚非欧的波斯帝国（当时世界上最大的帝国）两次入侵希腊。第一次主要是雅典在马拉松平原以少胜多击退了波斯大军——据说一位士兵从马拉松跑了约42公里回雅典报告胜利的消息，然后力竭而死。这就是今天"\textbf{马拉松赛跑}"的由来。

第二次规模更大：波斯国王薛西斯率领数十万大军（古希腊历史学家希罗多德说有两百多万，这个数字显然是夸张的）入侵，斯巴达国王列奥尼达率领300名斯巴达勇士在\textbf{温泉关}死守三天，全部战死。随后雅典人放弃了他们的城市（被波斯烧毁），全部转移到船上，在\textbf{萨拉米斯海战}中以少胜多击溃了波斯舰队。

希波战争的结局是希腊的胜利。如果波斯赢了——当时整个西方文明的走向可能完全改写。所以历史学家称希波战争为"\textbf{西方文明的生死时刻}"。

\subsection*{四、伯里克利时代与伯罗奔尼撒战争}

希波战争后，雅典进入了黄金时代（约公元前461—前429年），以当时的领袖\textbf{伯里克利}命名。\textbf{帕特农神庙}——那座至今仍然矗立在雅典卫城上的大理石神殿——就是伯里克利时代建的。哲学、戏剧、历史学都在这个时期达到了第一个高峰：苏格拉底在广场上和年轻人辩论，索福克勒斯写出了悲剧《俄狄浦斯王》，希罗多德被西塞罗称为"历史学之父"。

但这个黄金时代没能持续。雅典利用它在希波战争中获得的领导地位，把原来的防御同盟（提洛同盟）逐渐变成了自己的帝国——强迫盟邦交钱、驻军、接受雅典的统治。斯巴达领导的伯罗奔尼撒同盟对此无法容忍。公元前431年，\textbf{伯罗奔尼撒战争}爆发——两个最强大的希腊城邦互相厮杀，断断续续打了27年。最终雅典战败，整个希腊元气大伤。

\shihai{修昔底德陷阱——一个穿越两千五百年的警告}{
  雅典历史学家修昔底德亲身经历了伯罗奔尼撒战争，并写下了一部同名的历史巨著。他在书中提出了一个著名的判断——这场战争"不可避免的真正原因是雅典势力的增长和由此引起的斯巴达的恐惧"。

  这句话后来被浓缩成了一个政治学术语：\textbf{修昔底德陷阱}——当一个崛起中的大国挑战现有霸主时，战争往往不可避免。这个概念被现代国际关系学者反复引用，尤其是在讨论中美关系时——美国是当今世界的"霸主"，而中国是一个正在"崛起"的大国，两国的关系会不会陷入修昔底德陷阱？

  但修昔底德本人的态度其实更复杂。他写这部历史的目的是什么？他说的不是"战争不可避免所以算了"，而是希望通过他的记录，让\textbf{后人能够避免重蹈覆辙}。他的原话是："如果那些想要了解过去发生的事件……的人认为我的著作还有益处，我就满足了。"——历史不是用来预测命运的，是用来提供警告的。
}

\subsection*{五、亚历山大：把希腊文明带向世界}

希腊各城邦在伯罗奔尼撒战争中互相消耗得差不多了，北方的马其顿王国趁势崛起。马其顿国王腓力二世征服了希腊各城邦，但他的儿子\textbf{亚历山大}——亚里士多德的学生——做的远不止于此。

公元前334年，22岁的亚历山大率领一支约4万人的军队东征波斯。在接下来的十年里，他征服了小亚细亚、叙利亚、埃及、巴比伦、波斯本土，一直打到了印度河流域。他建立了一个横跨欧亚非三大洲的帝国——人类历史上第一个试图把东西方统一在同一个政权下的尝试。

亚历山大33岁在巴比伦病逝。他的帝国几乎立刻就分裂了（被他的将军们瓜分），但他留下了一个影响深远的遗产——\textbf{希腊化时代}。他把希腊的语言、哲学、艺术和城市生活方式带到了东方，在埃及建了亚历山大港（后来成为古代世界最大的学术中心），在各地建了数十座以自己名字命名的城市。不同文明之间的交流因此大大加速——这为后来罗马帝国的"普世化"铺垫了基础。

\subsection*{六、技术与文化}

\textbf{哲学：}
\begin{itemize}
  \item \textbf{苏格拉底}：不写书，只在雅典广场上和年轻人问答辩论。他的方法是不断追问，迫使对话者发现自己逻辑中的矛盾——即"\textbf{产婆术}"。他说"我知道我一无所知"——这不是谦虚，而是一种真正的哲学态度：承认无知是追求真知的起点。
  \item \textbf{柏拉图}（苏格拉底的学生）：创立了雅典学园（Academy），写了《理想国》，试图勾画一个由"哲学王"统治的完美城邦。英语academy（学院）这个词就来自他的学校。
  \item \textbf{亚里士多德}（柏拉图的学生）：亚历山大的老师，可能是古代世界最博学的人。他研究过的领域包括逻辑学、物理学、生物学、政治学、伦理学、诗学、修辞学——几乎涵盖了当时人类知识的所有门类。他的著作在中世纪被欧洲大学奉为教科书长达数百年。
\end{itemize}

\textbf{文学与戏剧：}
\begin{itemize}
  \item \textbf{荷马史诗}：《伊利亚特》和《奥德赛》是古希腊最早的文学巨著，口头流传数百年后才被写定。特洛伊战争、阿喀琉斯、奥德修斯的归乡之旅——这些故事至今是好莱坞电影和文学创作的素材库。
  \item \textbf{希腊悲剧}：埃斯库罗斯、索福克勒斯、欧里庇得斯三大悲剧作家的作品，探讨命运、自由意志、人性——这些主题到今天仍然是文学的核心母题。
\end{itemize}

\textbf{科学与数学：}
\begin{itemize}
  \item \textbf{毕达哥拉斯}（约前570—前495）：证明了"直角三角形斜边的平方等于两直角边的平方和"——我们今天还在数学课上学的勾股定理，在西方就叫毕达哥拉斯定理。
  \item \textbf{欧几里得}（约前300年）：著有《几何原本》，系统总结了此前的几何学知识。这本书是西方历史上发行量仅次于《圣经》的著作——直到20世纪初，它仍然是全世界学校的几何教材。
  \item \textbf{阿基米德}（约前287—前212）：希腊化时代最伟大的科学家。浮力定律（阿基米德定律）、杠杆原理、"给我一个支点，我能撬动地球"——都出自他。他还设计了各种巧妙的战争器械来保卫他的家乡叙拉古城。
\end{itemize}

\textbf{艺术与建筑：}
\begin{itemize}
  \item 帕特农神庙：多立克柱式的巅峰代表。它看起来四四方方，但实际上整个建筑里几乎\textbf{没有一条纯粹的直线}——地基微微上凸，柱子微微向内倾斜，柱间距微微不均匀调整。这些微妙的曲线矫正了人类视觉的错觉，使建筑从远处看起来完美挺直。
  \item 希腊雕塑：从早期的僵硬站姿（受埃及影响）发展到古典时期对\textbf{人体自然动态和理想化美感}的精确把握——米隆的《掷铁饼者》和波利克莱托斯的《持矛者》是典型代表。
\end{itemize}

\subsection*{七、本课要点回顾}

\begin{center}
\begin{tabular}{|c|c|p{7cm}|}
\hline
\textbf{时期} & \textbf{年代} & \textbf{核心事件} \\
\hline
城邦时代 & 前8世纪—前5世纪 & 斯巴达（军国）与雅典（民主）两种模式；梭伦改革、克里斯提尼改革 \\
\hline
希波战争 & 前499—前449年 & 马拉松（前490）、温泉关+萨拉米斯（前480）；希腊胜利 \\
\hline
黄金时代 & 前461—前429年 & 伯里克利；帕特农神庙；苏格拉底；戏剧繁荣 \\
\hline
伯罗奔尼撒战争 & 前431—前404年 & 雅典 vs 斯巴达；雅典战败；整个希腊元气大伤 \\
\hline
亚历山大东征 & 前334—前323年 & 征服波斯至印度；希腊化时代开始 \\
\hline
\end{tabular}
\end{center}

\vspace{0.3cm}
{\large\textcolor{orange!80!black}{\textbf{思考与讨论}}}
\begin{enumerate}
  \item 雅典的直接民主和今天我们国家的民主形式有什么不同？如果让你在"直接民主"和"代议制民主"之间做选择，你会选哪个？为什么？
  \item 雅典民主投票处死了苏格拉底——"多数人的决定"一定是对的吗？如果多数人投票决定剥夺少数人的基本权利，这个决定合法吗？
  \item 亚历山大征服了半个世界，却33岁就死了——他的帝国几乎没有在他身后存活。一个靠个人能力维持的帝国为什么难以持久？
  \item 从苏格拉底到亚里士多德，三代师生传承构成了西方哲学的基础。对比同时期中国春秋战国的百家争鸣——你能找出哪些相似和不同？
\end{enumerate}

\subsection*{要点总结}
\begin{itemize}
  \item 古希腊不是大一统帝国，而是由众多独立的城邦组成——其中斯巴达（军国）和雅典（民主）代表了两种极端模式
  \item 雅典民主是直接民主的早期实验，但仅限于成年男性公民——妇女、奴隶和外邦人被排除在外
  \item 苏格拉底之死揭示了"多数决"的危险：民主需要宪法和人权保护作为底线
  \item 希波战争的胜利保住了希腊文明的独立；伯罗奔尼撒战争则耗尽了希腊各城邦的元气
  \item 亚历山大东征开启了"希腊化时代"，将希腊文化传播到东方
\end{itemize}

\subsection*{讨论题}
\begin{enumerate}
  \item 斯巴达和雅典代表了两种完全不同的生活方式——一个极端强调集体，一个强调个人参与。你觉得今天的中国社会更接近哪一种？你自己更愿意生活在哪一种社会里？
  \item 雅典用抽签选官，现代民主用选举选官——抽签更"公平"（每个人都有机会），选举更"高效"（选到有能力的人）。你觉得这两种方式哪一种更好？有没有可能结合两者？
\end{enumerate}

\vspace{0.3cm}
\textcolor{blue!70!black}{\textbf{延伸阅读}}
\begin{itemize}
  \item 《写给孩子的希腊神话》——了解希腊文化的精神底色
  \item 纪录片《古希腊的瑰宝》——BBC出品，三集，从艺术角度解读希腊文明
\end{itemize}

\newpage

\section{\textcolor{orange!80!black}{第三课\quad 古罗马：从城邦到帝国，法律与征服的千年}}

当希腊城邦在伯罗奔尼撒战争中耗尽元气的时候，意大利半岛上一个不起眼的小城邦正在悄然崛起——\textbf{罗马}。

根据传说，罗马城建于公元前753年，建城者是罗慕路斯——一个被母狼养大的弃婴。这个传说是编的，但有一点说对了：\textbf{罗马从一开始就是一个充满野性的"狼性"文明}。

\subsection*{一、从王国到共和国：贵族与平民的博弈}

罗马早期是一个王国，但罗马人对国王的统治极其不满。公元前509年，罗马人驱逐了最后一任国王，宣布从此再也不设国王——\textbf{共和国}（Res Publica，拉丁语意为"公共事务"）诞生了。这一年，恰好也是雅典民主改革的前夜。东西方几乎在同一时间，以不同的方式，开启了"反对个人专制"的政治实验。

但罗马共和国不等于民主。真正的权力掌握在\textbf{元老院}——一个由贵族世家组成的小圈子——手中。国家最重要的行政官员（\textbf{执政官}，每年选举两人，互相制衡）几乎都来自几家世袭的贵族家族。

平民当然不满。公元前5世纪到前3世纪，罗马的平民展开了长达两百年的斗争——不是暴力革命（至少最初不是），而是不断\textbf{退出}：当国家面临外敌入侵时，平民集体拒绝服兵役——"你不给我们权利，我们就不替你打仗"。贵族一次次被迫妥协。最终，平民获得了选举\textbf{保民官}的权利——保民官有一项特殊的权力：\textbf{否决}。如果他认为元老院的决议损害了平民的利益，只要说一声"Veto"（拉丁语"我禁止"），元老院的决议就作废。

\shihai{罗马的"混合政体"——执政官 + 元老院 + 保民官}{
  罗马共和国的政治制度是一套精妙的权力制衡系统，后来对美国的建国者们产生了直接的影响：

  \begin{itemize}
    \item \textbf{执政官（代表王权）}——最高行政长官，两人共同执政，互相牵制。任期只有一年，不得连任。你想做什么大事？来得及吗？
    \item \textbf{元老院（代表贵族）}——三百人左右的精英议事机构，掌控财政、外交和重大决策。
    \item \textbf{保民官（代表平民）}——由平民选举产生，有权否决元老院的决议。
  \end{itemize}

  这三者的关系就像三个齿轮互相咬合：执政官想要大权独揽，元老院和保民官会联手阻止；元老院想要为所欲为，保民官会说Veto；保民官想要发动平民乱来，元老院和执政官会联手镇压。

  这套制度的核心精神是：\textbf{不信任任何一个单独的权力持有者。}罗马人设计的不是一个"好皇帝"制度——他们设计的是一个"没有皇帝"的制度。权力不是集中在一个人手里靠他的"英明"来保障，而是分散在不同势力之间靠\textbf{相互制衡}来维持稳定。
}

\subsection*{二、从共和国到帝国：民主为什么"自杀"了？}

罗马共和国在政治制衡上堪称精巧，在军事征服上更是恐怖。它先征服了意大利半岛，然后在地中海和迦太基（北非的一个海上贸易强国）进行了三次惨烈的\textbf{布匿战争}，最终灭亡迦太基。到公元前2世纪，罗马已经控制了大半个地中海。

但军事上的巨大成功给共和制度带来了致命的压力——\textbf{疆域太大，共和体制撑不住了}。

具体来说：战争给罗马带回了海量的奴隶和财富，但这些财富大部分流入了元老贵族的口袋。退伍的士兵回到家，发现自己的田地已经被大庄园吞并——罗马的平民阶层大规模破产，变成了涌入城市、靠国家救济度日的无产贫民。罗马军队最初是由有产公民组成的——你有地、买得起武器，你才有资格当兵。但公民普遍贫困化以后，兵源越来越紧缺。于是一个危险的改革出现了：将领们开始自己出钱招募穷人当兵，士兵的效忠对象从"共和国"变成了"给他们发军饷的将军"。军人私人化——这是共和制的掘墓人。

\textbf{尤利乌斯·凯撒}就是利用这种军阀力量上台的。他凭借在高卢（今法国）征战积累的军事威望和军队忠诚，于公元前49年渡过卢比孔河，率军进入罗马——这在共和时期是不可想象的背叛行为。凯撒后来被任命为"终身独裁官"，实际上已经等同于皇帝。公元前44年，一群元老在元老院将他刺杀——他们想拯救共和制，但已经太晚了，凯撒死了，共和制也死了。

凯撒的养子\textbf{屋大维}（后来的奥古斯都）在随后的内战中胜出。公元前27年，他玩了一个聪明的政治把戏：他向元老院宣布——"我把权力交还给你们，恢复共和。"元老院当然不敢要——谁敢要？他们反而恳求屋大维继续执政，并授予他"\textbf{奥古斯都}"（意为"神圣崇高者"）的称号。屋大维保留了共和国的一切机构——元老院照常开会，执政官照常选举——但所有实权都在他一个人手里。罗马\textbf{帝国}正式诞生，尽管它的统治者——\textbf{皇帝}——在很长时间里仍然自称不过是"第一公民"（Princeps）。

\shihai{罗马法——西方法律体系的源头}{
  罗马人对世界最持久的影响，可能不是疆域、不是建筑、不是凯撒——而是\textbf{法律}。

  罗马法是古代世界最系统、最精细的法律体系。它的核心精神可以概括为几点：
  \begin{itemize}
    \item \textbf{法律面前人人平等}（至少在自由民之间是平等的）
    \item \textbf{未经审判不得定罪}（这是现代"无罪推定"原则的萌芽）
    \item \textbf{私有财产不可侵犯}
    \item \textbf{契约必须遵守}（合同精神的基础）
  \end{itemize}

  公元6世纪，东罗马（拜占庭）皇帝查士丁尼组织法学家把千年来积累的罗马法汇编成《查士丁尼民法大全》——这是一部把散落在浩瀚法令中的罗马法整合成一体的超级工程。后来这部法典在中世纪晚期被重新发现，成为欧洲大陆各国制定民法典的基础。直到今天，\textbf{大陆法系}（以法国、德国为代表，世界上大多数国家包括中国都沿用这一体系）的直接源头就是罗马法。你日常生活中遇到的合同、继承、婚姻财产问题——处理这些问题的那套法律逻辑，很多都可以追溯到两千多年前的罗马法学家。
}

\subsection*{三、基督教的兴起与罗马帝国的分裂}

在罗马帝国的初期，遥远的犹太行省诞生了一个新的宗教——\textbf{基督教}。它的创始人是一个名叫耶稣的犹太木匠，他在大约30岁的时候开始传播他的教义——爱你的邻人、宽恕你的敌人、穷人将在天国得到安慰。耶稣被罗马总督钉死在十字架上（约公元30年），但他的追随者们——以保罗为代表——把基督教从一个小小的犹太教分支，发展成了罗马帝国境内最活跃的新宗教。

罗马帝国早期对基督教的态度是\textbf{迫害}——基督徒不拜罗马皇帝为神，这在罗马人看来等于叛国。但迫害没有消灭基督教，反而让它在底层民众中蔓延得更快。公元313年，君士坦丁大帝颁布"米兰敕令"，承认基督教的合法地位。到公元380年，基督教正式成为罗马帝国的国教。

与此同时，罗马帝国本身正走向分裂。公元395年，罗马帝国正式分为东西两部分：\textbf{西罗马帝国}（首都罗马）和\textbf{东罗马帝国}（首都君士坦丁堡，即拜占庭）。公元476年，西罗马帝国被日耳曼雇佣军首领废黜了最后一位皇帝，西罗马帝国灭亡。欧洲从此进入了\textbf{中世纪}。

东罗马帝国（拜占庭）则继续存在了近一千年——直到1453年被奥斯曼土耳其人攻克君士坦丁堡。

\subsection*{四、技术与文化}

\textbf{建筑与工程：}
\begin{itemize}
  \item 罗马人发明了\textbf{混凝土}（火山灰+石灰+碎石），这使他们能够建造前所未有的巨大穹顶和拱门。罗马万神殿的穹顶直径43.3米，直到19世纪之前一直是世界上最大的穹顶。
  \item \textbf{罗马大道}：条条大路通罗马——这句话不是夸张。罗马人修建了超过8万公里的石铺道路网络，从苏格兰边境一直延伸到波斯湾。这些道路的建设标准极高——多层碎石地基+排水沟+石铺路面，有些至今还能使用。
  \item \textbf{引水渠}：把几十公里外的山泉水通过高架水渠引入城市，供公共浴场、喷泉和居民使用。罗马的引水渠每天向城市输送约10亿公升的水——一个惊人的工程学成就。
\end{itemize}

\textbf{语言与文学：}
\begin{itemize}
  \item \textbf{拉丁语}是罗马的官方语言。罗马帝国灭亡后，拉丁语在各地逐渐演化为法语、西班牙语、意大利语、葡萄牙语、罗马尼亚语——这些统称"罗曼语族"。"罗曼"这个词本身就来自"罗马"。
  \item 古罗马文学：维吉尔的史诗《埃涅阿斯纪》（罗马版的"建国神话"）、西塞罗的演说和哲学著作、凯撒的《高卢战记》（他自己写的战争回忆录）——拉丁文学的黄金时代。
\end{itemize}

\subsection*{五、本课要点回顾}

\begin{center}
\begin{tabular}{|c|c|p{7cm}|}
\hline
\textbf{时期} & \textbf{年代} & \textbf{核心事件} \\
\hline
王政时代 & 前753—前509年 & 传说中罗慕路斯建城；罗马为王政 \\
\hline
共和国 & 前509—前27年 & 元老院+执政官+保民官制衡；平民与贵族斗争；布匿战争灭迦太基；凯撒独裁→遇刺 \\
\hline
帝国建立 & 前27年 & 屋大维称奥古斯都；保留共和外壳的帝制 \\
\hline
帝国全盛 & 1—2世纪 & 地中海成为"罗马的内湖"；罗马法体系形成 \\
\hline
基督教兴起 & 1—4世纪 & 耶稣传教→受难；保罗传播；313年米兰敕令合法化；380年成国教 \\
\hline
帝国分裂 & 395年 & 分为东西罗马 \\
\hline
西罗马灭亡 & 476年 & 被日耳曼人灭亡；欧洲进入中世纪 \\
\hline
\end{tabular}
\end{center}

\vspace{0.3cm}
{\large\textcolor{orange!80!black}{\textbf{思考与讨论}}}
\begin{enumerate}
  \item 罗马从共和走向帝制，是因为哪个关键环节出了问题？如果让你给当时的罗马设计一套防止军人私人化的制度，你会怎么做？
  \item 罗马法至今影响着世界大多数国家的法律体系——为什么一套两千年多前的法律可以"活"到今天？什么样的制度设计才具有跨时代的生命力？
  \item 基督教从一个被迫害的小教派变成罗马帝国的国教——这个转变是如何发生的？一个被排斥的弱势群体最终成为主流，这给你什么启发？
  \item 罗马人修建了惊人的道路和引水渠，中国的秦始皇也修了驰道和灵渠——两个同时代的大帝国不约而同地选择了大规模基础设施建设，这是为什么？
\end{enumerate}

\subsection*{要点总结}
\begin{itemize}
  \item 罗马共和国的混合政体（执政官+元老院+保民官）是权力制衡的早期典范——不信任任何一个单独的权力持有者
  \item 罗马从共和走向帝制的原因：领土扩张导致军队私人化、平民破产、军阀崛起
  \item 罗马法是西方法律体系的源头——法律面前人人平等、无罪推定、私有财产不可侵犯等原则影响至今
  \item 基督教从被迫害的小教派发展为罗马国教，改变了整个西方文明的精神底色
  \item 西罗马帝国灭亡后，东罗马（拜占庭）继续存在了近一千年
\end{itemize}

\subsection*{讨论题}
\begin{enumerate}
  \item 罗马共和国的政体设计核心是"不信任任何单一权力"——你觉得这种设计逻辑和中国的治理传统有什么根本不同？两种思路各有什么优劣？
  \item 罗马法的一些原则（如"契约必须遵守"）至今仍然有效——是什么样的特质让一套两千多年前的法律能"活"到今天？法律的生命力来自它的原则（抽象性）还是来自它的条文（具体性）？
\end{enumerate}

\vspace{0.3cm}
\textcolor{blue!70!black}{\textbf{延伸阅读}}
\begin{itemize}
  \item 纪录片《罗马：帝国无止境》——BBC出品，讲述罗马从城邦到帝国的全部历程
  \item 《罗马人的故事》——盐野七生著（少儿版），故事性极强的入门读物
\end{itemize}

\newpage

\section{\textcolor{orange!80!black}{第四课\quad 中世纪的欧洲与亚洲：在'黑暗'中酝酿新生}}

西罗马帝国在公元476年灭亡后，欧洲进入了\textbf{中世纪}——这个名称本身是后来文艺复兴时期的学者取的，意思是中间的"黑暗"时期。但用"黑暗"来概括整整一千年显然过于简单了——中世纪有它的混乱和倒退，也有它的创新和积累。

\subsection*{一、西欧封建制度：国王、领主与骑士}

西罗马帝国的中央集权崩溃之后，西欧变成了一个个互相争斗的小王国。没有哪个国王有能力建立起像罗马那样强大的官僚体系和常备军。于是，一种新的社会结构——\textbf{封建制度}——应运而生。

简单来说：国王把土地（采邑）分封给大贵族；大贵族再分封给小贵族（骑士）；骑士向封主宣誓效忠，每年提供一定天数的军事服务。农民（\textbf{农奴}）被绑定在土地上，没有自由迁徙的权利，终身为领主劳作。这种"土地换忠诚"的体系，在欧洲延续了数百年。

你可能会觉得这很像我们中国古代史里讲的分封制。确实有相似之处——但区别也很关键（我们在第一课的史海拾遗里已经比较过了：周制靠血缘，欧制靠契约）。这里补充一个实用的对比：欧洲封建制度中有一个著名的原则——"\textbf{我的附庸的附庸不是我的附庸}"。意思是国王不能直接指挥骑士——只能通过中间的大贵族。这导致国王的权力非常有限，很长一段时间里，法国国王直接控制的领地只有巴黎周边一小块，比好几个大贵族的领地都小。

正是这种\textbf{王权的弱小}，为后来西欧发展出了不同于东方的社会结构埋下了伏笔——贵族有和国王抗衡的实力，城市有和贵族谈判的资本，各种势力谁也吃不掉谁，只好坐下来谈条件、签合同、定法律。这培养了西欧社会特有的\textbf{契约精神}和\textbf{多元权力结构}。

\subsection*{二、基督教会的权力}

在中世纪的西欧，\textbf{基督教会}是唯一横跨各国边界的统一组织。教会的权力大到什么程度？它可以废黜国王——1077年，神圣罗马帝国皇帝亨利四世在风雪中赤脚站在卡诺莎城堡外三天三夜，向教皇格里高利七世忏悔认罪，史称"\textbf{卡诺莎之辱}"。

修道院在中世纪扮演了一个特殊的角色：它是知识的保存者。在识字率极低的时代，只有修道院的修士们还能读、能写、能抄书。我们今天能读到古希腊罗马的大量文献，很大程度上是因为中世纪的修士们一本一本、一页一页地把它们抄了下来。

\subsection*{三、阿拉伯帝国：当欧洲在'黑暗'中时，谁举着知识的火炬？}

当西欧在封建割据中缓慢前行的时候，东方一个新兴的文明正在迅速崛起——\textbf{阿拉伯帝国}。

公元7世纪，穆罕默德在阿拉伯半岛创立了\textbf{伊斯兰教}。他的追随者在不到一百年的时间里，建立了一个西起西班牙、东到印度河流域的庞大帝国。阿拉伯帝国的首都巴格达，在9世纪是世界上最大的城市之一，人口超过100万。当西欧的修道院还在小心翼翼地抄写残破的羊皮纸手稿时，巴格达的\textbf{智慧宫}里，来自波斯、印度、希腊、埃及的学者正在用阿拉伯语进行大规模的翻译和研究工作。

\shihai{百年翻译运动——阿拉伯帝国如何保存了古希腊的智慧}{
  从8世纪到10世纪，阿拉伯帝国的哈里发（统治者）们做了一件对世界文明影响深远的事：他们派出使者去拜占庭、波斯甚至印度搜集古代典籍，然后在巴格达组织学者大规模翻译。柏拉图、亚里士多德、欧几里得、托勒密、希波克拉底——这些古希腊哲人的著作，在基督教化的欧洲要么已经失传、要么被教会封存，但在阿拉伯帝国被完整地翻译成了阿拉伯语。

  阿拉伯学者的贡献绝不仅仅是"保存"——他们还做了大量创造性的工作：
  \begin{itemize}
    \item \textbf{代数}：花拉子密在9世纪写成了《代数学》，algebra（代数）这个英文单词就来自阿拉伯语al-jabr。他系统阐述了一次和二次方程的解法。
    \item \textbf{医学}：伊本·西拿（阿维森纳）写的《医典》被欧洲医学院用作教材直到17世纪——整整六百多年。
    \item \textbf{化学}：阿拉伯学者发展了蒸馏、结晶、过滤等实验技术。alcohol（酒精）、alkali（碱）、alchemy（炼金术）这些英文词汇都来自阿拉伯语。
  \end{itemize}

  后来这些著作通过西班牙和西西里岛传回了西欧——12世纪的一场"拉丁翻译运动"把阿拉伯语的希腊经典重新翻译回拉丁语。欧洲人这才重新发现了自己的"老祖宗"——而保存和传递这个火种的，是阿拉伯人。文明的传承从来不是某一条直线的延续，而是一个\textbf{多文明接力}的过程。
}

\subsection*{四、东方的同时代：从日本大化改新到宋元繁荣}

把视线转回东方。中世纪的一千年里，亚洲发生的事情同样波澜壮阔：

\textbf{日本：大化改新（645年）}——日本派出遣唐使到中国学习唐朝的制度和文化，回国后进行了全方位的"唐化"改革：废除贵族私有土地制、建立中央集权的律令制度、仿照长安建造了平城京（奈良）和平安京（京都）。这是世界史上最成功的"国家借鉴"案例之一——日本几乎把唐朝的制度、文字、佛教、建筑、历法全盘搬回了家。

\textbf{朝鲜半岛：新罗统一}——7世纪中后期，新罗在唐朝的援助下统一了朝鲜半岛大部分地区，同样吸收了唐朝的制度和文化，建立起中央集权的官僚国家。

\textbf{中国：隋唐→宋元}——这期间中国经历的是我们中国古代史课里讲过的隋唐盛世和宋元商业革命，对亚洲各国产生了强大的辐射作用。东亚的"汉字文化圈"和儒家文明圈，就是在这个时期定型的。

\subsection*{五、中世纪晚期的变革}

从11世纪开始，西欧的面貌开始发生变化：

\textbf{城市的复兴}。在罗马帝国崩溃后的几个世纪里，西欧的城市几乎消失了。但从11世纪起，随着农业技术的改进（重犁、三圃制）和商业的恢复，城市重新出现。意大利的威尼斯、热那亚，北欧的吕\rube{贝}{bèi}克、布鲁日，成为国际贸易的重要节点。

\textbf{大学的诞生}。最早的大学——意大利的博洛尼亚大学（1088年）和法国的巴黎大学（约1150年）——就是在这些新兴城市里出现的。中世纪大学和修道院不同：修道院是"关起门来修行"的地方，大学是"各色人等聚集辩论"的地方。法学、医学、神学、哲学在大学里被系统地研究和教授。现代大学的学位制度（学士、硕士、博士）就是中世纪遗留下来的。

\textbf{十字军东征}（1096—1291年）。西欧基督教世界发动了多次军事远征，名义上是收复被穆斯林占领的圣地耶路撒冷。十字军在军事上最终失败了，但在文化上产生了意外的影响：西欧的士兵和朝圣者第一次大规模接触到了比他们更先进的阿拉伯和拜占庭文明——精美的纺织品、香料、糖、亚里士多德的哲学——这些"舶来品"大大刺激了西欧人对东方财富和知识的渴望。间接地，十字军东征为后来的大航海时代埋下了动机。

\subsection*{六、技术与文化}

\textbf{农业技术：}
\begin{itemize}
  \item \textbf{重犁}：装上轮子和铁制犁头，比罗马时代轻便的扒犁更适合北欧厚重的黏土。配合新的挽具（套在牛肩而非牛脖子上），效率成倍提升。
  \item \textbf{三圃制}：把耕地分成三块，轮换种植冬小麦、春作物和休耕——土地利用率从二分之一提高到了三分之二，粮食产量大幅增加。
\end{itemize}

\textbf{建筑：}
\begin{itemize}
  \item \textbf{哥特式大教堂}：12世纪起源于法国，尖拱、飞扶壁、彩色玻璃窗和极高的中殿让教堂内部充满神秘的光影——巴黎圣母院、沙特尔大教堂都是哥特式建筑的巅峰代表。在技术上，飞扶壁把屋顶的侧推力传递到外部独立的扶\rube{垛}{duǒ}上，墙体可以大幅削减——从而可以在墙上开出巨大的窗户。
\end{itemize}

\textbf{文化：}
\begin{itemize}
  \item 日本的《源氏物语》（紫式部著，约1008年）被认为是世界上第一部长篇小说。
  \item 阿拉伯的《一千零一夜》（即《天方夜谭》）收集了波斯、印度和阿拉伯世界的民间故事，是中世纪最著名的文学瑰宝之一。
\end{itemize}

\subsection*{七、本课要点回顾}

\begin{center}
\begin{tabular}{|c|c|p{7cm}|}
\hline
\textbf{地区} & \textbf{时期} & \textbf{核心发展} \\
\hline
西欧 & 5—15世纪 & 封建制度（土地换忠诚）；基督教权力巅峰；11世纪后城市+大学复兴；哥特式建筑 \\
\hline
阿拉伯 & 7—13世纪 & 伊斯兰教兴起；百年翻译运动保存古希腊经典；代数学、医学领先世界 \\
\hline
日本 & 7世纪— & 大化改新全面学唐；汉字文化圈形成 \\
\hline
中国 & 6—14世纪 & 隋唐盛世→宋元商业革命（详见《中国古代史》） \\
\hline
\end{tabular}
\end{center}

\vspace{0.3cm}
{\large\textcolor{orange!80!black}{\textbf{思考与讨论}}}
\begin{enumerate}
  \item 西欧封建制度和西周分封制有什么本质区别？（可以回顾《中国古代史》第一课的"史海拾遗"）不同的制度选择对后来东西方的发展道路产生了什么影响？
  \item 阿拉伯帝国在8—10世纪是世界科学的中心——一个起于沙漠的民族是如何在短时间内成为学术霸主的？
  \item 日本大化改新大规模复制唐朝制度——这种"全盘借鉴"的模式有什么好处和风险？
  \item 西欧在11世纪后城市开始复兴——城市为什么是文明进步的发动机？中国的宋朝同时期也有城市繁荣（《清明上河图》），这两种城市化有什么异同？
\end{enumerate}

\subsection*{要点总结}
\begin{itemize}
  \item 中世纪西欧的封建制度以"土地换忠诚"为核心——"我的附庸的附庸不是我的附庸"导致王权弱小
  \item 基督教会是中世纪西欧唯一横跨各国边界的统一组织，修道院保存了古典文化的火种
  \item 阿拉伯帝国在8—10世纪是世界科学的中心——百年翻译运动保存并发展了古希腊学术
  \item 11世纪后西欧城市复兴、大学诞生，为中世纪末期的变革积蓄了力量
  \item 日本大化改新是中国制度向外输出的典型案例——几乎全盘复制了唐朝的制度文化
\end{itemize}

\subsection*{讨论题}
\begin{enumerate}
  \item "我的附庸的附庸不是我的附庸"——这种多层级的封建关系对后来欧洲的"契约精神"和"多元权力结构"有促进作用。你觉得"权力分散"和"权力集中"哪一种更有利于长期发展？历史上有定论吗？
  \item 阿拉伯帝国在鼎盛时期保存并发展了古希腊学术，然后通过西班牙传回西欧。你认为什么样的社会条件最有利于知识和学术的繁荣？是富裕、是开放、还是政治稳定？
\end{enumerate}

\vspace{0.3cm}
\textcolor{blue!70!black}{\textbf{延伸阅读}}
\begin{itemize}
  \item 纪录片《中世纪思潮》——以专题方式剖析中世纪的生活、信仰、学术
  \item 阿拉伯帝国相关：纪录片《伊斯兰：信仰帝国》（PBS出品）
\end{itemize}

\newpage

\section{\textcolor{orange!80!black}{第五课\quad 大航海时代：世界从分裂走向连接}}

15世纪，世界格局即将发生一场颠覆性的变革。而这场变革的起点，不是某位帝王的一声令下，而是一样不起眼的东西——\textbf{香料}。

今天的我们很难理解香料在中世纪的欧洲有多珍贵。没有冰箱的年代，保存肉类靠的是盐腌和香料腌制。胡椒、肉桂、丁香、豆\rube{蔻}{kòu}，来自遥远的印度和东南亚，经过波斯商人、阿拉伯商人、威尼斯商人一层一层的倒手，到了欧洲消费者手里，价格翻了几十倍甚至上百倍——胡椒被称为"黑色黄金"。而控制香料贸易陆路通道的关键枢纽——\textbf{君士坦丁堡}——在1453年被奥斯曼土耳其人攻陷。陆路被掐住了咽喉，欧洲人被迫把目光投向了海洋。

\subsection*{一、为什么是葡萄牙和西班牙？}

开启大航海时代的，不是当时欧洲最强的法国或英国，而是两个偏居伊比利亚半岛的小国家——\textbf{葡萄牙}和\textbf{西班牙}。

这不是偶然。第一，它们地理位置天然面向大西洋，向西向南都是未知的海洋；第二，它们刚刚完成了"收复失地运动"（把穆斯林政权赶出了伊比利亚半岛），积累了一批闲着没事做的骑士和冒险家；第三——也是最重要的——\textbf{亨利王子}（葡萄牙国王的弟弟）的远见。他在葡萄牙西南端的萨格里什建立了一所航海学校，召集了制图师、天文学家、造船工匠，集中攻关航海技术。这不是一个人的冒险，而是一个\textbf{国家级的项目}。

\subsection*{二、两条路线，一个目标：找到印度}

葡萄牙走的是\textbf{向东绕非洲}的路线。1488年，迪亚士绕过了非洲最南端的\textbf{好望角}。1498年，\textbf{达·伽马}成功绕过好望角抵达了印度的卡利卡特——香料之路打通了。当达·伽马返回里斯本时，他带回来的香料价值是他整个远航成本的60倍。

西班牙走了一条更冒险的路线——\textbf{向西横渡大西洋}。一个名叫\textbf{克里斯托弗·哥伦布}的意大利航海家，坚信地球是圆的（这在当时已经不是什么"异端邪说"——受过教育的人基本上都知道地球是圆的，争论的是地球有多"大"）。哥伦布认为向西航行很快就能抵达亚洲，但他严重低估了地球的周长。西班牙王室资助了他。

1492年8月3日，哥伦布率领三艘小船（最大的一艘也就二十多米长）从西班牙出发。两个多月后——10月12日——他们看到了陆地。哥伦布至死都相信自己抵达的是亚洲（所以他称当地居民为"Indians"——印度人）。实际上，他抵达的是今天巴哈马群岛中的一个小岛。他"发现"的，是一块欧洲人此前全然不知的大陆——\textbf{美洲}。

\shihai{哥伦布交换——一次改变地球生态的物种大搬家}{
  哥伦布"发现"新大陆（在美洲原住民看来，是欧洲人"被发现了"）引发了一场人类历史上规模最大的物种、疾病和文化的双向流动，历史学家称之为\textbf{"哥伦布交换"}（Columbian Exchange）。

  从旧大陆（欧亚非）到新大陆（美洲）：
  \begin{itemize}
    \item 小麦、水稻、甘蔗、咖啡
    \item 马、牛、猪、羊（美洲此前没有大型驯化动物）
    \item 天花、麻疹、流感——美洲原住民对这些旧大陆疾病毫无免疫力，据估计在1492年后的一个世纪里，美洲原住民人口减少了约90\%。这个数字极其惨痛，但它是客观历史的一部分。
  \end{itemize}

  从新大陆到旧大陆：
  \begin{itemize}
    \item 玉米、土豆、番薯、木薯——这些高产作物传入旧大陆后，极大地丰富了全球的食物供给。中国清代人口从一亿暴增到三亿，番薯和玉米在山区的广泛种植是一个重要原因。
    \item 番茄、辣椒、烟草、可可（巧克力）
  \end{itemize}

  你今天的餐桌上——土豆丝、番茄炒蛋、烤玉米、辣椒炒肉——这些食物在1492年之前的欧亚大陆是不存在的。它们全部来自美洲。\textbf{"哥伦布交换"把以前各自独立演化的两个生态系统首次连接在了一起}，人类从此进入了一个"全球生物共同体"的时代——有甜蜜，也有悲剧。
}

\subsection*{三、麦哲伦的环球航行}

1519年，葡萄牙航海家\textbf{麦哲伦}（为西班牙服务）率领五艘船从西班牙出发，向西横渡大西洋，绕过南美洲南端的麦哲伦海峡（他发现并命名的），横渡太平洋（他给这个"太平"的海洋起的名字，因为当时风平浪静），抵达菲律宾。在菲律宾，麦哲伦因卷入当地部落冲突被杀，没能活着回到欧洲。但他的船队继续航行，穿过印度洋，绕过好望角——1522年，仅剩的一艘船"维多利亚号"载着18名幸存者回到了西班牙。人类第一次用实践证明：\textbf{地球确实是圆的}，环绕一圈需要三年。

\subsection*{四、殖民与掠夺：大航海时代的另一面}

大航海不仅是冒险与发现的故事，也是暴力与掠夺的故事。西班牙征服者\textbf{科尔特斯}带着几百人就灭亡了拥有数百万人口的阿兹特克帝国（墨西哥），\textbf{皮萨罗}用更少的人灭亡了印加帝国（秘鲁）。西班牙人从美洲运回了天文数字的黄金和白银——据估计，16—17世纪西班牙从美洲掠取的白银超过1.6万吨。这些贵金属流入欧洲，造成了所谓的"\textbf{价格革命}"——白银泛滥导致物价飞涨，传统的封建经济体系受到巨大冲击，加速了资本主义在西欧的萌芽。

大航海还催生了人类历史上最黑暗的制度之一——\textbf{三角贸易}。欧洲的船只满载枪支、纺织品到西非换取奴隶；奴隶被塞进船舱横渡大西洋（"中程航线"，死在这段路上的人数以百万计）；到了美洲他们被卖给种植园主；船只再装满蔗糖、棉花、烟草返回欧洲。一个三角回路，每一个顶点都沾满了鲜血。

\shihai{郑和下西洋 vs 哥伦布——两种不同的"走向海洋"}{
  1405年—1433年，郑和七下西洋（详见《中国古代史》第七课），比哥伦布早了将近90年。郑和的宝船比哥伦布的旗舰大几十倍，船队规模之大更是当时的欧洲航海家无法想象的。但郑和死后，中国主动退出了海洋。

  而哥伦布的远航开启的是一个全新的时代——殖民扩张、全球贸易、物种交换。为什么明朝有最强的船却选择了放弃，而欧洲后来居上？

  几个关键因素的差异：
  \begin{itemize}
    \item \textbf{经济驱动力不同：}郑和下西洋是朝贡外交——花钱扬国威，不赚钱。哥伦布远航是商业投资——西班牙王室是要分利润的。前者靠财政拨款，断了就没了；后者形成了可持续的经济循环。
    \item \textbf{政治结构不同：}明朝是一个统一的陆地帝国，它的核心利益是守住北方的农耕区。海洋对它来说是"可有可无的外围"——就算放弃了，帝国仍然能稳定运转。而欧洲是多个竞争性小国，一个国家的国王如果不支持航海，人才和资本就会流到邻国去。这种多国竞争的压力让海上探险"停不下来"。
    \item \textbf{地理不同：}中国面向一片巨大的太平洋，向东走几万公里才到美洲。欧洲西面就是大西洋，哥伦布出发两个多月就"撞"上了美洲——地理的偶然为欧洲的崛起提供了便利。
  \end{itemize}

  这个对比没有简单的好坏之分——但它说明：\textbf{决定历史走向的，往往不是一个文明"聪明不聪明"，而是它面临的地理约束、经济激励和政治结构。}
}

\subsection*{五、技术与文化}

\textbf{航海技术：}
\begin{itemize}
  \item \textbf{卡拉维尔帆船}：葡萄牙人改良的轻快帆船，吃水浅、逆风也能走（利用三角帆），适合远洋探险。
  \item \textbf{罗盘和星盘}：指南针（从中国传入后改良）和天文导航工具，让远洋航行不再靠肉眼猜方向。
  \item \textbf{波特兰海图}：标注了海岸线、港口和罗盘方向的实用航海图——大航海时代的GPS。
\end{itemize}

\textbf{地理学：}
\begin{itemize}
  \item 托勒密的《地理学》在15世纪初被重新翻译成拉丁文——虽然它对地球周长的估算偏小（哥伦布正是被这个错误"鼓舞"而向西航行的），但它系统记录了经纬度的概念和当时已知的世界地理。
  \item 大航海时代的地理大发现彻底改变了欧洲人绘制世界地图的方式。1570年，奥特里乌斯出版了第一本现代意义上的世界地图集。
\end{itemize}

\subsection*{六、本课要点回顾}

\begin{center}
\begin{tabular}{|c|c|p{7cm}|}
\hline
\textbf{关键事件} & \textbf{年代} & \textbf{意义} \\
\hline
葡萄牙探航 & 15世纪 & 亨利王子航海学校；迪亚士发现好望角（1488）；达伽马到印度（1498） \\
\hline
哥伦布"发现"新大陆 & 1492年 & 开启旧大陆与新大陆的持续接触 \\
\hline
麦哲伦环球航行 & 1519—1522年 & 首次实证地球是圆的 \\
\hline
殖民扩张 & 16—17世纪 & 西班牙在美洲掠夺金银；三角贸易 \\
\hline
\end{tabular}
\end{center}

\vspace{0.3cm}
{\large\textcolor{orange!80!black}{\textbf{思考与讨论}}}
\begin{enumerate}
  \item 郑和下西洋和哥伦布远航（可在《中国古代史》第七课找到对比）——为什么技术领先的一方反而退出了海洋竞争？
  \item "哥伦布交换"带来了巨大的经济利益，也带来了人口灾难。面对这种"硬币的两面"，你怎么评价大航海时代的整体影响？
  \item 葡萄牙和西班牙这样的小国能够率先开启大航海时代——国家的大小和它的创新能力有关系吗？
  \item 三角贸易是大航海时代最黑暗的一页。为什么当时整个欧洲社会似乎都"默认"了这种罪恶的制度？经济利益能够为不道德的行为提供合理性吗？
\end{enumerate}

\subsection*{要点总结}
\begin{itemize}
  \item 大航海时代的直接推动力是香料贸易和奥斯曼帝国切断陆路通道——欧洲被迫寻找海上新航路
  \item 葡萄牙和西班牙率先开启大航海，得益于地理位置、王室支持和航海技术的集中攻关（亨利王子的航海学校）
  \item 哥伦布"发现"美洲引发了"哥伦布交换"——旧大陆和新大陆之间的物种、疾病和文化的双向大迁移
  \item 大航海带来了殖民掠夺、三角贸易和奴隶制——这是现代全球化光明与黑暗交织的起点
  \item 郑和下西洋与哥伦布远航的对比揭示：不同的经济驱动力和政治结构导致了截然不同的海洋战略
\end{itemize}

\subsection*{讨论题}
\begin{enumerate}
  \item 有人说哥伦布是"英雄"（连接了两个世界），也有人说他是"恶魔"（开启了殖民掠夺）。你觉得一个人物的历史评价应该以他的"意图"为准还是以"后果"为准？我们今天对历史人物的评价标准在变化吗？
  \item 郑和和哥伦布的对比说明：即使技术领先，如果缺乏持续的经济动机，领先也会被放弃。你还能举出其他"技术领先但没有转化为持续优势"的例子吗？
\end{enumerate}

\vspace{0.3cm}
\textcolor{blue!70!black}{\textbf{延伸阅读}}
\begin{itemize}
  \item 纪录片《大航海时代》——BBC出品，追踪哥伦布、麦哲伦等航海家的真实航程
  \item 贾雷德·戴蒙德《枪炮、病菌与钢铁》——解答"为什么是欧洲征服了美洲而不是反过来"这个根本问题（少儿版可先从简写本入手）
\end{itemize}

\newpage

\section{\textcolor{orange!80!black}{第六课\quad 文艺复兴与宗教改革：'人'的重新发现}}

如果用一个词概括中世纪末期到近代早期欧洲的精神史，这个词就是——\textbf{人的发现}。

在大约两个世纪的时间里（14—16世纪），欧洲经历了文艺复兴和宗教改革两场巨变。前者把"人"从上帝的影子下解放了出来，后者把信仰从教会的垄断中解放了出来。

\subsection*{一、文艺复兴：为什么从意大利开始？}

\textbf{文艺复兴}（Renaissance，法语"再生"之意）发源于14世纪的意大利，不是偶然的。意大利有三个得天独厚的条件：

第一，\textbf{钱}。意大利北部城市——佛罗伦萨、威尼斯、米兰——是中世纪欧洲最富裕的地区，商业和银行业极其发达。美第奇家族——佛罗伦萨的实际统治者——就是靠银行业起家的超级富豪。他们资助了一大批艺术家和学者——米开朗基罗、达·芬奇、波提切利都得到过美第奇家族的赞助。有了钱，艺术和学术才能繁荣——"\textbf{美第奇效应}"后来成了"富豪赞助艺术和科学"的代名词。

第二，\textbf{古希腊罗马的遗产就在脚下}。意大利遍地是罗马帝国的废墟——你随便挖个地下室都可能碰到古罗马的雕像残片。而1453年君士坦丁堡陷落后，大批拜占庭学者带着古希腊文献逃到了意大利，带来了西欧几乎已经丢失的希腊古典文化。\textbf{回到古代源头去}——这是文艺复兴的核心口号。

第三，\textbf{城市自治传统}。意大利城邦（不像法国和英国那样在王权统治下）由富商和银行家统治，思想氛围相对宽松自由。

\textbf{文艺复兴的"灵魂"}可以用一个词概括——\textbf{人文主义}。在此之前的中世纪欧洲，一切价值标准都与上帝相关：人生的意义是死后上天堂，艺术的功能是赞美上帝，学术的职责是解释《圣经》。而人文主义者说：\textbf{人本身就有价值，人的理性、情感、创造力是值得赞美的，人可以在现世——而不只是死后——过上有意义的生活}。

这不是"反对宗教"（文艺复兴时期几乎所有大艺术家都画过宗教题材），而是\textbf{把焦点从神转移到人}。达·芬奇画的《最后的晚餐》虽然还是基督教题材，但他画的耶稣和门徒不再是神圣不可亲近的符号，而是一个个有血有肉、有情感反应的真实人物。

\shihai{文艺复兴三杰与他们的'跨界天才'}{
  文艺复兴时期最让人惊叹的，也许是那些\textbf{不满足于只做一件事的人}。

  \textbf{列奥纳多·达·芬奇}（1452—1519）：最著名的身份是画家（《蒙娜丽莎》《最后的晚餐》），但他同时也是雕塑家、建筑师、音乐家、数学家、工程师、解剖学家、地质学家、植物学家。他设计了飞行器、坦克、潜水服的草图，解剖了几十具尸体来研究人体结构——他的笔记本里画满了肌肉、骨骼和胚胎的精细素描。这些工程和科学设想绝大多数在他生前没有实现，但现代工程师按照他的设计图造出了能飞的滑翔机——达·芬奇领先了他的时代约四百年。

  \textbf{米开朗基罗}（1475—1564）：雕塑（《大卫》）、绘画（西斯廷教堂天顶画《创世纪》）、建筑（圣彼得大教堂穹顶）、诗歌——四个领域都达到了同时代最高水平。他的大卫雕像高达4米多，从一块被其他雕塑家"雕坏了"的弃石中"拯救"出来——他对人体结构的理解精确到了每一块肌肉的起伏。

  \textbf{拉斐尔}（1483—1520）：只活了37岁，但他的《雅典学院》把古希腊以来的所有伟大哲学家和科学家画在了同一个想象中的大厅里——柏拉图指天（理念世界）、亚里士多德指地（经验世界）——这是他献给人文主义的理想图卷。

  为什么文艺复兴能产生这么多"全才"？一个重要原因是\textbf{知识尚未分化}——当时没有"你是学文科还是理科"这种分界线，一个好奇的人可以去研究所有他感兴趣的东西。今天我们的知识体系被分成了越来越细的专业，这提高了研究效率，但也让"达·芬奇式"的跨界天才越来越稀少了。
}

\subsection*{二、宗教改革：一个人挑战整个教会}

如果说文艺复兴是从\textbf{文化}上松动教会的垄断，宗教改革则是从\textbf{信仰}上给了它沉重一击。

1517年10月31日（这一天后来成了宗教改革的纪念日），德意志地区一位大学神学教授\textbf{马丁·路德}把一份长长的清单——《\textbf{九十五条论纲}}——钉在了维滕堡教堂的大门上。这是一篇学术辩论邀请，但他猛烈抨击的对象是教会最赚钱也最腐败的一项做法：\textbf{赎罪券}。

教会宣称：你犯了罪，可以通过购买赎罪券来缩短你的灵魂在炼狱中受苦的时间——"钱币叮当一声落入银箱，灵魂就能从炼狱飞升天堂。"马丁·路德认为这是彻头彻尾的欺诈。他的核心理由是：\textbf{人和上帝之间的关系是个人的、直接的，不需要教会做中介。}一个人能否得救，取决于他对上帝的信仰（"因信称义"），而不是他买了多少赎罪券、做了多少善功。

这个主张在今天听起来可能只是"神学争论"，但放在16世纪，它是一枚\textbf{核弹}。因为教会的全部权力——它的财富、它的政治影响、它对人们生活的控制——都建立在一个前提上：\textbf{教会是人与上帝之间必不可少的中介}。如果人人可以直接和上帝对话，那还需要教会干什么？

马丁·路德拒绝收回他的主张，被教皇开除教籍并宣布为异端。按照惯例他应该被烧死——但他没死。因为他得到了萨克森选帝侯腓特烈三世的保护。腓特烈保护路德不是因为他认同路德的神学——他是出于\textbf{政治动机}：削弱教皇对德意志各邦的控制。这个细节很重要——\textbf{宗教改革得以成功，不只是神学上的正确，更是政治上碰巧有人需要它}。

\shihai{古腾堡印刷术——宗教改革的技术加速器}{
  马丁·路德不是第一个批评教会腐败的人。在他之前一个多世纪，波希米亚的扬·胡斯就提出了类似的抗议——他在1415年被烧死在火刑柱上。

  为什么路德没有被烧死？除了政治庇护之外，还有一个关键的技术因素：\textbf{印刷术}。

  1450年左右，德意志人\textbf{约翰内斯·古腾堡}（比毕\rube{昇}{shēng}晚约四百年）发明了金属活字印刷术和手摇印刷机。这套技术让书籍的复制成本大幅降低——在印刷术出现之前，一本书要靠修士手抄几个月甚至几年；有了印刷机，一天就能印出数百页。

  马丁·路德的《九十五条论纲》在贴在教堂门上之后，不到两周就被印刷成小册子传遍了德意志，不到一个月传遍了整个西欧。到他去世时，他的著作已经卖出了数十万册。如果没有印刷术，路德很可能就是另一个被烧死的胡斯；有了印刷术，他的思想在教会还来不及组织有效镇压之前就已经深入人心了。

  这印证了一个历史规律：\textbf{技术往往是思想革命的加速器。}新思想需要新的传播工具——没有印刷术，就没有宗教改革；没有互联网，很多当代的社会运动也达不到今天的规模和速度。
}

\subsection*{三、宗教改革的后续影响}

路德之后，宗教改革在欧洲迅速蔓延。瑞士的\textbf{加尔文}建立了比路德更激进的新教体系，英格兰国王\textbf{亨利八世}因为想离婚而与教皇闹翻（不是为了什么崇高的神学信念，就是因为他想换一个能生儿子的老婆），干脆宣布英国教会脱离罗马教皇的控制，自任英国教会的最高领袖——这就是英国\textbf{国教}的起源。

宗教改革导致了一系列血腥的宗教战争——德意志的三十年战争（1618—1648）可能是欧洲历史上死人最多的战争之一，整个中欧被打成了废墟。\textbf{威斯特伐利亚和约}（1648年）最终确立了一条现代国际关系的基础原则：\textbf{一个国家的主权不容他国以宗教为借口进行干涉}。这标志着以宗教统合欧洲的中世纪秩序的终结，也是现代\textbf{民族国家}体系的起点。

\subsection*{四、技术与文化}

\textbf{艺术：}
\begin{itemize}
  \item \textbf{透视法}：文艺复兴时期发展起来的\textbf{线性透视法}（近大远小、平行线汇聚于消失点）是绘画史上的一次视知觉革命。它让二维的画布可以逼真地呈现三维空间——在马萨乔的《圣三位一体》里，你感觉墙壁真的被凿开了一个洞。
  \item \textbf{油画技术}：凡·艾克兄弟（15世纪尼德兰画家）改良了油画颜料和技法，使得画面色彩更丰富、层次更细腻——油画的"油画性"是从这里开始的。
\end{itemize}

\textbf{科学革命的先声：}
\begin{itemize}
  \item \textbf{哥白尼}（1473—1543）：提出了日心说——不是太阳围着地球转，而是地球围着太阳转。这是一个颠覆性的观念，因为它不仅在科学上挑战了托勒密的地心说，更在\textbf{世界观}上撼动了"人类是宇宙中心"的信念。
  \item \textbf{伽利略}（1564—1642）：用望远镜观察天空，看到了月球表面的环形山、木星的卫星、金星的相位变化。他用\textbf{实验}来检验理论——从一个斜面上滚下小球，测量时间和距离的数学关系——这是近代科学方法的奠基。伽利略因此被教会审判并被迫"认错"——科学和宗教的冲突，在他身上体现得最尖锐。
\end{itemize}

\textbf{文学：}
\begin{itemize}
  \item \textbf{但丁《神曲》}（14世纪初）：用意大利俗语而不是拉丁语写的长诗，是人文主义的文学先声。
  \item \textbf{莎士比亚}（1564—1616）：可能是英语文学史上最伟大的作家。《哈姆雷特》里的"To be, or not to be"——对生与死的沉思——是人文主义精神在文学中的极致表达：人开始认真思考自己的存在意义，而不再只是等着神告诉他答案。
  \item \textbf{塞万提斯《堂吉诃德》}（1605—1615）：被公认为现代小说的开山之作——一个读骑士小说读疯了的老人骑着瘦马去行侠仗义，处处碰壁却不肯放弃他的理想。
\end{itemize}

\subsection*{五、本课要点回顾}

\begin{center}
\begin{tabular}{|c|c|p{7cm}|}
\hline
\textbf{事件} & \textbf{年代} & \textbf{核心内容} \\
\hline
文艺复兴 & 14—16世纪 & 发源于意大利（美第奇赞助）；人文主义：从以神为中心转向以人为中心 \\
\hline
三杰 & 15—16世纪 & 达·芬奇、米开朗基罗、拉斐尔；透视法、人体解剖、理想美 \\
\hline
宗教改革 & 1517年— & 马丁·路德《九十五条论纲》抨击赎罪券；"因信称义"；古腾堡印刷术加速传播 \\
\hline
科学革命 & 16—17世纪 & 哥白尼日心说；伽利略的实验方法 \\
\hline
文学 & 14—17世纪 & 但丁《神曲》、莎士比亚戏剧、塞万提斯《堂吉诃德》 \\
\hline
\end{tabular}
\end{center}

\vspace{0.3cm}
{\large\textcolor{orange!80!black}{\textbf{思考与讨论}}}
\begin{enumerate}
  \item 文艺复兴的"人文主义"和中世纪的"以上帝为中心"有什么本质区别？你今天的生活中，哪些方面还能看到人文主义的影响？
  \item 古腾堡印刷术对宗教改革的成功至关重要——一项技术发明如何改变了一场思想运动？今天有没有类似的"技术加速思想传播"的例子？
  \item 达·芬奇横跨艺术和科学的多个领域——今天为什么很难再看到这样的"全才"？知识分科太细是好事还是坏事？
  \item 马丁·路德和伽利略都挑战了教会的权威，但他们的结局截然不同（路德善终，伽利略受审）。是什么导致了这种差异？
\end{enumerate}

\vspace{0.3cm}
\textcolor{blue!70!black}{\textbf{延伸阅读}}
\begin{itemize}
  \item 纪录片《美第奇家族：现代艺术的缔造者》
  \item 《达·芬奇传》——沃尔特·艾萨克森著（少儿版），了解这位史上最著名跨界天才的传奇一生
  \item 莎士比亚戏剧——推荐先看改编的儿童版或动画版，如《王子复仇记》（即《哈姆雷特》）、《罗密欧与朱丽叶》
\end{itemize}

\subsection*{要点总结}
\begin{itemize}
  \item 文艺复兴发源于意大利北部，因为有商业财富（美第奇家族赞助）、古典遗产和城市自治传统——核心精神是"人文主义"，把焦点从神转移到人
  \item 达·芬奇、米开朗基罗、拉斐尔代表了文艺复兴艺术和科学的高峰，透视法和人体解剖是技术的核心突破
  \item 马丁·路德的宗教改革挑战了教会的权威——"因信称义"使人和上帝的关系不再需要教会中介
  \item 古腾堡印刷术是宗教改革的技术加速器——新思想的传播速度决定了谁能在镇压到来前赢得人心
  \item 文艺复兴和宗教改革共同松动了中世纪的教会垄断，为科学革命和启蒙运动开辟了道路
\end{itemize}

\subsection*{讨论题}
\begin{enumerate}
  \item 文艺复兴为什么能产生达·芬奇这样的"跨界天才"？今天的教育体系鼓励"专精"还是"跨界"？你觉得自己以后想成为"专才"还是"通才"？
  \item 马丁·路德说人可以绕过教会直接和上帝对话——这在当时为什么是革命性的？你能想到今天社会中有没有类似的"去中介化"运动（去掉中间人，让个体直接连接）？
\end{enumerate}

\section*{讨论提示}

\subsection*{第一课（四大文明古国）讨论题提示}
\begin{itemize}
  \item 关于"地理决定论 vs 其他因素"：可以思考——地理是文明发展的起点，但不是全部。同样在大河流域，中国文明延续至今，而巴比伦和印度河流域的早期文明中断了。地理解释了"为什么在那里开始"，但不能解释"为什么延续或中断"。
  \item 关于"技术先进的文明为何脆弱"：可以思考——哈拉帕的高度城市化可能让它更依赖复杂的系统（贸易、粮食供应），系统越复杂就越脆弱——一个环节断了整个链条就停了。小而简单的社会反而更有韧性。
\end{itemize}

\subsection*{第二课（古希腊）讨论题提示}
\begin{itemize}
  \item 关于"斯巴达 vs 雅典"：可以思考——斯巴达的集体主义保障了安全和稳定但扼杀了个人自由；雅典的个人自由催生了文化繁荣但导致了内部分裂和内耗。今天的你可能同时需要两者——在安全的环境下拥有自由的空间。
  \item 关于"抽签 vs 选举"：可以思考——抽签确保平等但可能选出无能者；选举确保能力但可能让有钱有势者垄断。古希腊的抽签和现代的选举可以结合——比如陪审团用抽签，行政官用选举。
\end{itemize}

\subsection*{第三课（古罗马）讨论题提示}
\begin{itemize}
  \item 关于"制衡 vs 集中"：可以思考——制衡防止权力滥用但可能降低决策效率；集中提高效率但可能走向专制。中国治理传统更侧重效率和集中，西方治理传统更侧重制衡——两种模式各有利弊，适应不同的历史条件和社会需求。
  \item 关于"法律的生命力"：可以思考——罗马法的原则之所以能存活至今，恰恰是因为它足够抽象（如"契约必须遵守"）——抽象的原则可以适应不同的时代条件。过于具体的条文会随着时代变迁而失效。
\end{itemize}

\subsection*{第四课（中世纪）讨论题提示}
\begin{itemize}
  \item 关于"权力分散 vs 权力集中"：可以思考——中世纪欧洲的权力分散（封建制）意外地催生了多元权力结构和契约传统；而东方的权力集中带来了高效治理但也可能导致僵化。两者各有利弊，关键看适应什么样的时代条件。
  \item 关于"学术繁荣的条件"：可以思考——阿拉伯帝国的学术繁荣发生在政治统一但文化开放的时代（巴格达智慧宫吸引了各族学者）。对比同时期的欧洲，分裂但竞争激烈也是创新的温床。开放、富裕和竞争可能比稳定更有利于学术。
\end{itemize}

\subsection*{第五课（大航海时代）讨论题提示}
\begin{itemize}
  \item 关于"意图 vs 后果"：可以思考——哥伦布个人可能并无恶意（他至死认为自己到了亚洲），但他的航行带来的后果是灾难性的（原住民人口锐减90\%）。今天我们在评判历史人物时，越来越倾向于用"后果"而非"意图"——这本身是一种价值观的变化。
  \item 关于"技术领先为何未能持续"：可以思考——除了郑和之外，还有英国在工业革命早期对蒸汽机的领先（后来被德国和美国追赶）。技术领先不能自动转化为持续优势——需要制度、资本和市场共同支撑。
\end{itemize}

\subsection*{第六课（文艺复兴与宗教改革）讨论题提示}
\begin{itemize}
  \item 关于"专才 vs 通才"：可以思考——达·芬奇时代知识总量少，一个人可以跨多个领域。今天知识爆炸，任何人都不可能什么都懂。但这不意味着跨界不重要——真正的大创新往往发生在学科交界处。你需要"一根足够深的钉子（专业）+ 足够宽的视野（跨界）"。
  \item 关于"去中介化"：可以思考——路德去掉的是教会这个宗教"中介"；今天互联网去掉的是媒体、零售、出行等行业的中介。"去掉中介"听起来解放了个体，但也可能让个体失去保护——没有人过滤信息真伪、没有人保障交易安全。
\end{itemize}

\vspace{0.5cm}
\begin{center}
\textcolor{blue!70!black}{\large —— 六课走过四千年世界文明，从金字塔到环球航行，历史是全人类共同的故事。下一站：近现代史！ ——}
\end{center}

\end{document}
